% Part: normal-modal-logic
% Chapter: filtrations
% Section: finite

\documentclass[../../../include/open-logic-section]{subfiles}

\begin{document}

\olfileid{nml}{fil}{fin}

\olsection{ફિલ્ટ્રેશનો સાંત છે}

ઉપ-!!{formula}s હેઠળ બંધ એવા કોઈપણ ગણ~$\Gamma$ માટે ફિલ્ટ્રેશન
વ્યાખ્યાયિત કર્યું છે. વ્યાખ્યા પોતે ફિલ્ટ્રેશન સાંત હશે એની ખાતરી
આપતી નથી. ખરેખર, $\Gamma$ અનંત હોય (દાખલા તરીકે, બધાં
!!{formula}sનો ગણ હોય), તો ફિલ્ટ્રેશન પણ અનંત હોઈ શકે. પરંતુ
$\Gamma$ સાંત હોય (દાખલા તરીકે, આપેલા !!{formula}~$!A$નાં
ઉપ-!!{formula}sનો ગણ હોય), તો~$\Gamma$ દ્વારા કોઈપણ ફિલ્ટ્રેશન
પણ સાંત હોય છે.

\begin{prop}\ollabel{prop:filt-are-finite}
  જો $\Gamma$ સાંત હોય, તો નિદર્શન $\mModel{M}$નું $\Gamma$ દ્વારા
  કોઈપણ ફિલ્ટ્રેશન $\mModel{M^*}$ પણ સાંત છે.
\end{prop}

\begin{proof}
  $W^*$નું કદ એ સમકક્ષતા સંબંધ~$\equiv$ હેઠળના જુદા વર્ગો~$[w]$ની
  સંખ્યા છે. આવા વર્ગનાં કોઈપણ બે જગત $u$, $v$—એટલે કે એવાં
  $u$ અને $v$ કે~$u \equiv v$—$\Gamma$નાં બધાં
  !!{formula}s~$!A$ પર સહમત હોય છે: $!A \in \Gamma$ હોય તો $!A$
  બંને $u$ અને $v$ પર $!A$ સત્ય હોય અથવા બંને પર અસત્ય હોય. તેથી દરેક
  વર્ગ~$[w]$ને~$\Gamma$નો એક ઉપગણ અનુરૂપ છે:~$[w]$નાં જગતો પર
  સત્ય હોય એવાં બધાં $!A \in \Gamma$નો ગણ. બે જુદા વર્ગ
  $[u]$ અને $[v]$ એ~$\Gamma$ના એ જ ઉપગણને અનુરૂપ હોઈ શકે નહિ. કારણ કે જો
  $u$ પર સત્ય !!{formula}sનો ગણ અને $v$ પર સત્ય !!{formula}sનો
  ગણ એક જ હોય, તો $u$ અને $v$ એ~$\Gamma$નાં બધાં સૂત્રો પર
  સહમત હોય, એટલે $u \equiv v$. પરંતુ ત્યારે $[u] = [v]$.
  આમ $W^*$માંથી $\Pow{\Gamma}$માં એક !!a{injective} વિધેય છે અને
  તેથી $\card{W^*} \le \card{\Pow{\Gamma}}$. જો $\Gamma$માં $n$
  વાક્યો હોય, તો $W^*$નું કદ વધુમાં વધુ~$2^n$ છે.
\end{proof}

\end{document}
